\scanpage{162}
\mspage{36}
% archive sheet number: 161

Le triple $(W,\unsure{E'|C'},F')$ est un triple de Quillen clos
(correspondant à une ``closed model category'') de
Quillen, \inserted{du moins} si on suppose $M$ stable par $\underrightarrow{\lim}$ finies et
$\underleftarrow{\lim}$ finies, et que $W$ \deleted{satisfasse} est modérément saturé).

% margin note written sideways
\margin{Il en résulte le \fbox{Cor. 4} \deleted{Pour que $(W,C,F)$ soit un \unsure{triple}, de Quillen clos,} \inserted{Soit $W$ un loc. mod. \unsure{saturé}} \deleted{il faut que $C$ clos à g. et $F$ clos à droite, $(C,TF\overset{\mathrm{déf}}=C_*)$} et \unsure{$(TC\overset{\mathrm{déf}}=C\cap W$} et $(C,TF),(TC,F)$ deux \inserted{pré-}couples de Quillen tels que $TC\subset C\cap W$, $TF\subset F\cap W$. Les cond. suivantes sont équiv. a) $TC=C\cap W$ b) $TF=F\cap W$ c) $W=TF\circ TC$ d) $(W,C,\unsure{B|F})$ est un \deleted{\unsure{triple}} de Quillen clos.}

\uline{Démonstration} Les relations (4.3.1) impliquent
\[
\widetilde{TC}\subset\widetilde C\quad\text{et}\quad\overline{TF}\subset\overline F
\]
i.e.
\[
TC'\subset C',\quad T\unsure{P|F}'\subset F',
\]
i.e. (4.3.1) pour les nouveaux couples de
Quillen. \inserted{La relation} 4.3.2. $W\supset TC,TF$, implique
\[
W\supset\widetilde{TC},\overline{TF},\quad\text{i.e.}\quad W\supset TC',TF',
\]
\deleted{\unsure{en}} en vertu du cor. 2 à la prop. 4.2 et du
fait que $W$ est stable par facteurs directs ;
\deleted{comme} donc on a (4.3.2) pour les nouveaux
couples. Comme par hypothèse on a
\[
W=\overline{TF}\circ\widetilde{TC}
\]
et ceci s'écrit
\[
W=TF'\circ TC',
\]
donc les nouveaux couples satisfont aux $\unsure{=\mathrm{s}|\Rightarrow}$ conditions,
qui cette fois, les couples étant clos, prennent
la forme forte (4.3.6).

Prouvons que le triple $(W,C',F')$ satisfait
aux axiomes de Quillen d'une \deleted{structure} \inserted{triple} de
Quillen clos. On a bien, par 4.3.6.
\[
\left\{\begin{aligned}
&C'=\inserted{(TF')^*={}}(F'\cap W)^*,\quad F'=(TC')_{\unsure{*|1}}=(C'\cap W)_*\\
&W=TF'\circ TC'=C'_*\circ {F'}^{*}
\end{aligned}\right.
\]
qui correspond à la définition des ``catégories de
modèles de Quillen \uline{clos}''. Les axiomes
M1, \deleted{\unsure{(pro.}} M2, M3, M4 de Quillen sont \unsure{réalisés|rédigés}
de façon triviale, \deleted{M5 \unsure{par hypoth}} (M3 et M4 puisque
